\documentclass[conference]{IEEEtran}

\usepackage[T1]{fontenc}
\usepackage[utf8]{inputenc}
\usepackage{amsmath,amssymb,amsfonts}
\usepackage{booktabs}
\usepackage{graphicx}
\usepackage{multirow}
\usepackage{cite}
\usepackage{url}
\usepackage{hyperref}
\hypersetup{hidelinks}

\begin{document}

	itle{Dual-Channel FFT + KAN Deepfake Detection:\\
A Reproducible Experimental Report on DeepDetect-2025}

\author{\IEEEauthorblockN{Satvik Pathak}
\IEEEauthorblockA{Department/Institution Name \\
City, Country \\
email@domain.com}}

\maketitle

\begin{abstract}
This paper reports the exact experimental pipeline implemented in \texttt{bestfileyet007.ipynb} for binary deepfake detection (Real=0, Fake=1). The method converts grayscale spatial images to a dual-channel frequency representation and classifies samples with a convolutional stem followed by Kolmogorov--Arnold Network (KAN) layers. On a balanced DeepDetect-2025 subset of 103{,}986 images (51{,}993 real, 51{,}993 fake), with a deterministic 70/15/15 split, the run reports test loss 0.3683, accuracy 0.8324, precision 0.8222, recall 0.8483, F1-score 0.8351, and ROC-AUC 0.9172. This manuscript includes only formulas, settings, and outcomes directly present in the notebook implementation and outputs.
\end{abstract}

\begin{IEEEkeywords}
Deepfake Detection, Frequency Domain, FFT, Phase Spectrum, Magnitude Spectrum, Kolmogorov--Arnold Networks.
\end{IEEEkeywords}

\section{Scope and Experimental Contract}
This manuscript documents one executed pipeline and one reported run from \texttt{bestfileyet007.ipynb}. The report is constrained to verified items only:
\begin{itemize}
    \item Input/output task: binary classification (Real vs Fake).
    \item Data source used in code: DeepDetect-2025 (downloaded via \texttt{kagglehub}).
    \item Model used in code: dual-channel FFT preprocessing + \texttt{PhaseKAN}.
    \item Metrics reported from notebook outputs: accuracy, precision, recall, F1, ROC-AUC, confusion matrix, and classification report.
\end{itemize}
No external baseline, zero-shot, or state-of-the-art claim is made here.

\section{Dataset Construction and Split}
The notebook discovers image files and infers labels from directory names:
\begin{itemize}
    \item label 0: paths containing \texttt{real}
    \item label 1: paths containing \texttt{fake}/\texttt{ai}/\texttt{synthetic}, or generator keywords
\end{itemize}
A balanced subset is built by truncating both classes to equal size.

Observed manifest (printed by the notebook run):
\begin{itemize}
    \item Total: 103{,}986
    \item Real: 51{,}993
    \item Fake: 51{,}993
\end{itemize}

Let $N=103{,}986$. With ratios $(0.70,0.15,0.15)$ and integer rounding in code:
\begin{align}
N_{\text{train}} &= \lfloor 0.70N \rfloor = 72{,}790, \\
N_{\text{val}}   &= \lfloor 0.15N \rfloor = 15{,}597, \\
N_{\text{test}}  &= N - N_{\text{train}} - N_{\text{val}} = 15{,}599.
\end{align}

\section{Preprocessing and Frequency-Domain Formulation}
\subsection{Spatial preprocessing}
For each sample, the code performs:
\begin{enumerate}
    \item PIL load
    \item grayscale conversion (single channel)
    \item crop to $224\times224$ (random crop + horizontal flip for train, center crop for val/test)
    \item conversion to tensor in $[0,1]$
\end{enumerate}
No resize interpolation is used in the notebook pipeline.

\subsection{GPU dual-channel FFT extraction}
Given a batch $\mathbf{X}\in\mathbb{R}^{B\times1\times H\times W}$, the code computes
\begin{equation}
\mathbf{F}=\operatorname{fftshift}(\operatorname{fft2}(\mathbf{X}_{\text{squeeze}})).
\end{equation}

Two channels are then created:
\begin{align}
\mathbf{P} &= \angle(\mathbf{F}), \\
\mathbf{P}_{\text{norm}} &= \frac{\mathbf{P}+\pi}{2\pi}, \\
\mathbf{M} &= \log\left(|\mathbf{F}| + 10^{-8}\right), \\
\mathbf{M}_{\text{norm}} &= \frac{\mathbf{M}-\min(\mathbf{M})}{\max(\mathbf{M})-\min(\mathbf{M})+10^{-8}},
\end{align}
where min/max are computed per sample over spatial dimensions.
The final model input is
\begin{equation}
\mathbf{Z}=\operatorname{stack}[\mathbf{P}_{\text{norm}},\mathbf{M}_{\text{norm}}]\in\mathbb{R}^{B\times2\times224\times224}.
\end{equation}

\section{Model and Optimization Equations}
\subsection{Implemented architecture}
The notebook model is \texttt{PhaseKAN}:
\begin{itemize}
    \item Conv stem: $(2\!\to\!64,s=2)\to(64\!\to\!128,s=1)\to(128\!\to\!128,s=2)\to(128\!\to\!256,s=1)\to(256\!\to\!256,s=2)$ with BatchNorm and GELU.
    \item Adaptive average pool to a 256-d vector.
    \item KAN head with widths $256\to128\to64\to1$ (with LayerNorm and Dropout=0.3 between hidden layers).
\end{itemize}
Reported trainable parameters: 1{,}337{,}154.

\subsection{KAN block used in code}
For input vector $\mathbf{x}$, each implemented \texttt{KANLinear} block uses two additive paths:
\begin{equation}
\mathbf{y}=\mathbf{W}_s\,\phi(\mathbf{x})+\mathbf{b}_s + \mathbf{W}_b\,\operatorname{SiLU}(\mathbf{x})+\mathbf{b}_b.
\end{equation}
Here $\phi(\cdot)$ is a per-feature B-spline activation module with trainable coefficients.

\subsection{Loss and decision rule}
The notebook uses \texttt{BCEWithLogitsLoss}. For logit $z$ and label $y\in\{0,1\}$:
\begin{equation}
\mathcal{L}(z,y) = -\left[y\log\sigma(z) + (1-y)\log(1-\sigma(z))\right].
\end{equation}
Prediction rule used in code:
\begin{equation}
\hat{y}=\mathbb{1}[\sigma(z)\ge 0.5].
\end{equation}

\subsection{Training hyperparameters from code}
\begin{itemize}
    \item Epoch budget: 35
    \item Optimizer: AdamW, learning rate $10^{-3}$, weight decay $10^{-3}$
    \item Scheduler: CosineAnnealingLR ($T_{\max}=35$, $\eta_{\min}=10^{-6}$)
    \item Gradient clipping norm: 1.0
    \item Early stopping patience: 10
    \item Batch size: 96
    \item Mixed precision on CUDA: enabled
\end{itemize}

\section{Reported Results (Notebook Output)}
\subsection{Primary test metrics}
\begin{table}[!t]
\caption{Test-set metrics printed by the notebook.}
\label{tab:primary}
\centering
\begin{tabular}{lc}
	oprule
Metric & Value \\
\midrule
Loss & 0.3683 \\
Accuracy & 0.8324 (83.24\%) \\
Precision & 0.8222 \\
Recall & 0.8483 \\
F1-score & 0.8351 \\
ROC-AUC & 0.9172 \\
\bottomrule
\end{tabular}
\end{table}

\subsection{Class-wise and confusion matrix values}
The classification report printed by the notebook:
\begin{itemize}
    \item Real: precision 0.8433, recall 0.8165, F1 0.8297, support 7799
    \item Fake: precision 0.8222, recall 0.8483, F1 0.8351, support 7800
\end{itemize}

Confusion matrix (rows=true [Real, Fake], columns=pred [Real, Fake]):
\begin{equation}
\mathbf{C}=\begin{bmatrix}
6368 & 1431 \\
1183 & 6617
\end{bmatrix}.
\end{equation}

\subsection{Training trajectory values}
The logged run shows 35 completed epochs and restoration of the best validation-loss checkpoint. The smallest printed validation loss is 0.3552 (epoch 32), with validation accuracy 0.8418 at that epoch.

\section{Reproducibility Snapshot (Printed)}
\begin{itemize}
    \item Python: 3.12.12
    \item PyTorch: 2.10.0+cu128
    \item Device: CUDA
    \item Seed: 42
\end{itemize}

Notebook-generated files reported in output:
\begin{itemize}
    \item \texttt{dual\_channel\_fft\_ieee.png}
    \item \texttt{kan\_deepfake\_results.png}
    \item \texttt{runs/kan\_deepfake/best\_phase\_kan.pt}
    \item \texttt{runs/kan\_deepfake/training\_summary.json}
\end{itemize}

\section{Limitations (Observed Scope Only)}
This report covers one dataset construction pipeline and one run output from the notebook. It does not include external baseline comparison tables, multi-seed statistics, or cross-dataset zero-shot evaluation in the presented factual results.

\section{Conclusion}
The implemented dual-channel FFT + KAN pipeline in \texttt{bestfileyet007.ipynb} produces the reported test metrics on the DeepDetect-2025 balanced subset, with ROC-AUC 0.9172 and accuracy 83.24\%. All equations and numbers in this manuscript are directly aligned with the notebook code and printed outputs.

\begin{thebibliography}{00}
\bibitem{kan2024} Z. Liu \emph{et al.}, ``KAN: Kolmogorov--Arnold Networks,'' 2024.
\end{thebibliography}

\end{document}